\subsection{Організація і методика проведення педагогічного експерименту}\label{sec:3.2}

\subsubsection{Експериментальна база та учасники дослідження}

Педагогічний експеримент проводився на базі кафедри інформатики та прикладної математики Криворізького державного педагогічного університету (КДПУ) протягом 2020--2025~рр. Вибір експериментальної бази зумовлений тим, що КДПУ є провідним закладом вищої освіти Криворізького регіону з підготовки вчителів інформатики, де розроблення та впровадження імерсивних освітніх технологій є одним із пріоритетних напрямів наукових досліджень кафедри~\cite{Semerikov2022IERDesign}.

У дослідженні брали участь студенти спеціальності 014 Середня освіта (Інформатика) другого (магістерського) освітнього рівня, що навчалися на денній формі навчання. Загальна кількість учасників експерименту~-- 64 особи, з яких 34 студенти склали експериментальну групу (ЕГ) і 30 студентів~-- контрольну групу (КГ).

Формування ЕГ та КГ здійснювалося з урахуванням таких вимог:
\begin{itemize}
\item приблизно однаковий рівень успішності студентів за результатами попередніх семестрів;
\item однаковий обсяг навчальних дисциплін інформатичного блоку;
\item однакові умови матеріально-технічного забезпечення навчального процесу (доступ до комп'ютерних лабораторій, мобільних пристроїв для тестування WebAR, мережі Інтернет);
\item відсутність попереднього досвіду систематичного навчання розробки ІОР.
\end{itemize}

Експериментальна група навчалася за розробленою методикою з використанням авторського навчально-методичного забезпечення (навчальний посібник з WebAR, курс <<Синтетична реальність в освіті>> у LMS Moodle, лабораторний практикум, віртуальна хімічна лабораторія). Контрольна група навчалася за традиційною програмою підготовки, що передбачала вивчення окремих аспектів WebAR у межах дисциплін інформатичного блоку без системного та послідовного навчального забезпечення.

\subsubsection{Етапи педагогічного експерименту}

Дослідно-експериментальну роботу було організовано у три послідовні етапи: констатувальний (2020--2021), формувальний (2021--2024) та контрольний (2024--2025).

\textbf{I. Констатувальний етап (2020--2021)} мав на меті вивчення початкового стану сформованості вмінь розробки ІОР у студентів спеціальності 014 Середня освіта (Інформатика), визначення рівня їхньої обізнаності з технологіями WebAR/VR/XR, аналіз мотивації до вивчення імерсивних технологій, а також підтвердження статистичної еквівалентності ЕГ та КГ перед початком формувального етапу.

На констатувальному етапі було виконано:
\begin{enumerate}
\item проведено вхідне анкетування студентів щодо обізнаності з технологіями AR/VR/XR та досвіду їх використання (результати засвідчили, що понад 80\% студентів мали лише побіжне уявлення про доповнену реальність, переважно на рівні використання AR-фільтрів у соціальних мережах, і жоден не мав досвіду розробки WebAR додатків);
\item здійснено діагностику початкового рівня сформованості вмінь розробки ІОР за розробленим критеріальним апаратом (підрозділ~\ref{sec:3.1});
\item проведено порівняльний аналіз інструментів розробки WebAR (A-Frame, AR.js, Three.js, JSARToolKit, Babylon.js), результати якого опубліковано у \cite{Shepiliev2020WebAR}; цей аналіз дав змогу обґрунтувати вибір MindAR та Three.js як основних засобів навчання розробки WebAR;
\item розроблено та апробовано профорієнтаційні квести з використанням WebAR на базі КДПУ, що продемонстрували потенціал імерсивних технологій в освіті та підтвердили доцільність їх системного вивчення \cite{Shepiliev2021Career};
\item проведено статистичну перевірку еквівалентності ЕГ та КГ за критерієм $\chi^2$ Пірсона, що підтвердила відсутність статистично значущих відмінностей між групами ($p > 0{,}05$).
\end{enumerate}

\textbf{II. Формувальний етап (2021--2024)}~-- основний етап дослідно-експериментальної роботи, протягом якого в експериментальній групі було впроваджено розроблену методику навчання розробки ІОР з використанням авторського навчально-методичного забезпечення.

Мета формувального етапу~-- перевірити ефективність розробленої методики навчання, що спирається на структурно-функціональну модель підготовки (підрозділ~\ref{subsec:2.2}) та реалізує визначені педагогічні умови (підрозділ~\ref{subsec:2.3}).

Завдання формувального етапу:
\begin{enumerate}
\item забезпечити послідовне та системне навчання студентів ЕГ розробки WebAR додатків~-- від основ тривимірної графіки у браузері до інтеграції моделей машинного навчання;
\item реалізувати педагогічні умови навчання: психолого-педагогічні (створення мотивуючого навчального середовища), методичні (використання навчального посібника та відеоматеріалів), технологічні (забезпечення доступу до сучасних засобів розробки WebAR), організаційні (систематична практична діяльність протягом 18 тижнів);
\item здійснювати поточний моніторинг рівня сформованості вмінь студентів ЕГ за критеріями, визначеними у підрозділі~\ref{sec:3.1};
\item коригувати методику навчання на основі результатів поточного моніторингу.
\end{enumerate}

\textbf{III. Контрольний етап (2024--2025)} мав на меті підсумкове оцінювання рівня сформованості вмінь розробки ІОР у студентів ЕГ та КГ, порівняння результатів, статистичну обробку та інтерпретацію отриманих даних.

На контрольному етапі було виконано:
\begin{enumerate}
\item проведено підсумкову діагностику рівня сформованості вмінь розробки ІОР у студентів ЕГ та КГ за критеріальним апаратом (підрозділ~\ref{sec:3.1});
\item здійснено статистичну обробку результатів за критерієм $\chi^2$ Пірсона для перевірки значущості відмінностей між ЕГ та КГ;
\item проведено якісний аналіз динаміки змін за кожним критерієм;
\item сформульовано рекомендації щодо вдосконалення навчання розробки ІОР.
\end{enumerate}

\subsubsection{Навчально-методичне забезпечення формувального етапу}

Основу навчально-методичного забезпечення формувального етапу становлять три взаємопов'язані компоненти: навчальний посібник з WebAR, курс <<Синтетична реальність в освіті>> у LMS Moodle та комплект лабораторних завдань.

\textbf{Навчальний посібник <<Розробка WebAR-застосунків для освіти>>} є системним навчальним виданням, що послідовно розкриває всі аспекти розробки імерсивних освітніх ресурсів на основі технологій WebAR. Посібник структуровано у 10 розділів, що відповідають логіці поступового ускладнення навчального матеріалу:

\begin{enumerate}[label=\arabic*.]
\item \textit{<<Вступ до WebAR>>}~-- огляд технологій доповненої реальності у веб-браузері, їх переваги (крос-платформність, відсутність необхідності встановлення додатків) та обмеження; порівняння WebAR із нативними AR-додатками; знайомство з основними бібліотеками: MindAR, AR.js, Three.js, A-Frame, WebXR; визначення TensorFlow.js як засобу інтеграції машинного навчання.

\item \textit{<<Засоби розробки>>}~-- налаштування середовища розробки WebAR: основи JavaScript (змінні, типи даних, оператори, функції, масиви, об'єкти, DOM), встановлення та використання локальних веб-серверів, створення HTTPS-тунелів за допомогою ngrok для тестування на мобільних пристроях, налаштування віддаленого налагодження через Chrome DevTools, використання заздалегідь записаних відео для імітації камери під час розробки.

\item \textit{<<Як працює доповнена реальність у веб-браузері>>}~-- фундаментальні концепції WebAR: WebGL як основа тривимірного рендерингу у браузерах; бібліотека Three.js та її ключові компоненти (сцена, камера, рендерер); створення тривимірних об'єктів (геометрія, матеріали, текстури); налаштування освітлення; анімація; принцип накладання віртуальних 3D-об'єктів на відеопотік з камери мобільного пристрою.

\item \textit{<<Відстеження зображень (Image Tracking)>>}~-- бібліотека MindAR для відстеження природних зображень (Natural Feature Tracking): підготовка зображень, компіляція NFT маркерів, створення прив'язок тривимірних об'єктів до відстежуваних зображень; завантаження текстур, відео та 3D-моделей у форматі GLTF/GLB як AR-контенту; інтерактивність.

\item \textit{<<Відстеження обличчя (Face Tracking)>>}~-- модуль MindAR Face Tracking: робота з 468 опорними точками обличчя; прив'язка 3D-об'єктів до конкретних точок обличчя; створення окклюдерів для реалістичного накладання; маски та фільтри; перемикання камер; захоплення кадрів.

\item \textit{<<Відстеження довкілля (World Tracking)>>}~-- стандарт WebXR Device API: створення XR-сесій; компонент ARButton; управління контролерами; розміщення 3D-об'єктів у реальному просторі за допомогою hit testing; перевірка дотику (touch input).

\item \textit{<<Інтеграція машинного навчання>>}~-- трикрокова методика інтеграції моделей ML у WebAR: (крок~1) інтеграція стандартної моделі TensorFlow.js handpose; (крок~2) навчання власної моделі за допомогою Teachable Machine; (крок~3) інтеграція навченої моделі у WebAR додаток. Також бібліотека face-api.js для розпізнавання емоцій.

\item \textit{<<Інші технології WebAR>>}~-- A-Frame, model-viewer, комерційні платформи (8th Wall, Zappar).

\item \textit{<<Висновки>>}~-- узагальнення та систематизація розглянутих технологій.

\item \textit{<<Лабораторний практикум>>}~-- 10 варіантів творчих завдань з 3D-моделювання та AR.
\end{enumerate}

До навчально-методичного забезпечення також входить \textbf{кейс <<Віртуальна хімічна лабораторія у доповненій реальності>> (anchem-ar)}~-- інтеграційний проєкт, що демонструє реальне освітнє застосування WebAR. Кейс передбачає створення віртуальної лабораторії для виконання лабораторної роботи <<Якісні реакції хлорид-, бромід- та йодид-іонів>>, де кожний реактив (AgNO$_3$, Pb(NO$_3$)$_2$, HCl, NaBr, KI та ін.) співставляється з маркером доповненої реальності. При наближенні двох маркерів один до одного відбувається <<зливання>> реактивів~-- відтворюється відеозапис реального хімічного експерименту. Реалізація здійснюється засобами AR.js з використанням A-Frame та демонструє повний цикл розробки ІОР від педагогічної ідеї до програмної реалізації.

\textbf{Курс <<Синтетична реальність в освіті>> у LMS Moodle.} Навчальний курс розгорнуто на платформі moodle.kdpu.edu.ua для студентів другого (магістерського) освітнього рівня, другий семестр. Курс структуровано у тижневому форматі (18 тижневих модулів, лютий~-- червень) та включає такі тематичні блоки:

\begin{enumerate}[label=\arabic*.]
\item \textit{Тиждень~1:} Вступ до WebAR~-- ознайомлення з технологіями доповненої реальності у веб-браузері.
\item \textit{Тиждень~2:} Підготовка зображення для MindAR~-- налаштування середовища, компіляція NFT маркерів.
\item \textit{Тиждень~3:} Завантаження текстур~-- робота з текстурами у Three.js.
\item \textit{Тижні~4--5:} Завантаження 3D-моделей та відтворення відео~-- імпорт GLTF/GLB, мультимедійний AR-контент.
\item \textit{Тиждень~6:} Лабораторна робота~1 <<Візитівка у доповненій реальності>>.
\item \textit{Тижні~7--8:} Опорні точки обличчя та накладання маски~-- MindAR Face Tracking, 468 ландмарків.
\item \textit{Тиждень~9:} Перемикання камери та захоплення кадрів.
\item \textit{Тиждень~10:} Лабораторна робота~2 <<Примірка віртуальних аксесуарів>>.
\item \textit{Тижні~11--13:} WebXR Device API~-- ARButton, контролери, XR-сесії.
\item \textit{Тиждень~14:} Розміщення об'єктів~-- hit testing, розміщення на поверхнях.
\item \textit{Тижні~15--16:} WebAR та машинне навчання~-- TensorFlow.js handpose, face-api.js.
\item \textit{Тиждень~17:} Перевірка дотику~-- touch input для розміщення об'єктів.
\item \textit{Тиждень~18:} A-Frame, model-viewer, комерційні засоби.
\item \textit{Підсумкова атестація}~-- захист індивідуальних проєктів.
\end{enumerate}

Кожний тижневий модуль курсу включає: навчальний матеріал (розділ посібника у PDF та/або відеолекція); практичне завдання з програмування WebAR; тест або контрольні запитання; форум для обговорення та взаємодопомоги.

\textbf{Лабораторний практикум} складається з двох обов'язкових лабораторних робіт та комплексу з 10 варіантів творчих завдань.

\textit{Лабораторна робота~1 <<Візитівка у доповненій реальності>>} (тиждень~6) передбачає створення персональної AR-візитівки засобами MindAR Image Tracking, що включає 3D-об'єкти, текстури, анімацію та гіперпосилання.

\textit{Лабораторна робота~2 <<Примірка віртуальних аксесуарів>>} (тиждень~10) передбачає розробку WebAR додатку для <<примірки>> віртуальних аксесуарів за допомогою MindAR Face Tracking з використанням окклюдерів.

\textit{Творчі завдання} (10 варіантів) передбачають створення 3D-моделей реальних об'єктів з прив'язкою до маркерів AR: масштабована модель системи Земля--Місяць (варіант~0), храмовий комплекс Гізи (варіант~1), Ханойські вежі (варіант~2), Сатурн з кільцями (варіант~3), будівля Криворізької міської ради (варіант~4), планетарна модель атома літію (варіант~5), фізична симуляція пружного удару (варіант~6) та ін. Кожне завдання включає чотири кроки: побудова 3D-моделі $\to$ анімація $\to$ текстурування та освітлення $\to$ прив'язка до маркера AR.

\subsubsection{Методика кількісного оцінювання}

Кількісне оцінювання результатів експерименту здійснювалося на основі критеріального апарату (підрозділ~\ref{sec:3.1}), що включає 100-бальну рубрику з рівномірним розподілом між чотирма критеріями (по 25 балів). Для статистичної перевірки значущості відмінностей між ЕГ та КГ використано критерій $\chi^2$ Пірсона при рівні значущості $\alpha = 0{,}05$ та $k = 3$ ступенях свободи. Поточний моніторинг на формувальному етапі включав аналіз результатів тестування, оцінювання лабораторних робіт, аналіз активності у LMS Moodle та есе-рефлексій.

\subsubsection{Методика якісного аналізу транскриптів лекцій}\label{subsec:3.2.qual}

Для поглиблення аналізу навчального процесу на формувальному етапі було застосовано якісний аналіз транскриптів відеолекцій курсу <<Синтетична реальність в освіті>> як додатковий метод дослідження. Це відповідає принципам змішаного дослідницького дизайну~\cite{Creswell2018MixedMethods}, де кількісні та якісні дані збираються паралельно та інтегруються на етапі інтерпретації (конвергентний паралельний дизайн).

\textit{Обґрунтування.} Кількісні результати педагогічного експерименту (підрозділ~\ref{sec:3.1}) фіксують \textit{що} змінилося у рівнях сформованості вмінь, проте не розкривають \textit{як} саме навчальний процес забезпечив ці зміни. Аналіз педагогічного дискурсу дає змогу виявити конкретні механізми реалізації педагогічних умов, динаміку факторів іммерсивного навчання та стратегії викладання, що сприяли позитивній динаміці за кожним критерієм.

\textbf{Транскриптовий корпус.} Корпус складається з 15 файлів автоматичної транскрипції (ASR) відеолекцій, записаних протягом 18 тижнів курсу (лютий~-- червень). Характеристику корпусу подано у табл.~\ref{tab:corpus}.

\begin{table}[!h]
\centering
\caption{Характеристика транскриптового корпусу}
\label{tab:corpus}
\small
\begin{tabular}{|c|l|c|c|c|}
\hline
\textbf{Тижні} & \hfill \textbf{Тема} & \textbf{Сегм.} & \textbf{Хв.} & \textbf{Етап} \\
\hline
1 & Вступ до WebAR & 50 & 81 & ТА \\
\hline
2 & Підготовка зображення для MindAR & 58 & 87 & ТА \\
\hline
3 & Завантаження текстур & 63 & 80 & ТА \\
\hline
4--5 & Завантаження 3D-моделей & 105 & 153 & ТА \\
\hline
6--7 & Відтворення відео та лаб.~робота~1 & 108 & 155 & ТА \\
\hline
8 & Опорні точки обличчя (Face Tracking) & 51 & 80 & ПР \\
\hline
9 & Накладання маски на обличчя & 68 & 94 & ПР \\
\hline
10 & Перемикання камери та захоплення кадрів & 41 & 63 & ПР \\
\hline
11 & Початок роботи із WebXR & 58 & 83 & ПР \\
\hline
12 & Компонент ARButton & 50 & 91 & ПР \\
\hline
13 & Управління контролерами & 54 & 64 & ПР \\
\hline
14 & Розміщення об'єктів\textsuperscript{*} & --- & --- & --- \\
\hline
15 & Розпізнавання віку та статі (TensorFlow.js) & 56 & 89 & ІР \\
\hline
16--17 & Перевірка дотику та комерційні засоби & 132 & 139 & ІР \\
\hline
18 & 8th Wall та середовища розробки & 111 & 141 & ІР \\
\hline
\multicolumn{2}{|l|}{\textbf{Усього (14 файлів)}} & \textbf{1005} & \textbf{1399} & \\
\hline
\end{tabular}

\medskip
\footnotesize
\textit{Примітка.} Етапи: ТА~-- теоретико-аналітичний (тижні~1--6), ПР~-- практико-розвивальний (тижні~7--14), ІР~-- інтегративно-рефлексивний (тижні~15--18). \textsuperscript{*}Файл тижня~14 виключено: ASR-система розпізнала мовлення як англійську замість української.
\end{table}

\textbf{Попередня обробка.} Транскрипти у форматі VTT (Web Video Text Tracks) були оброблені програмно: вилучено часові мітки та службову розмітку; суміжні сегменти з перекриттям часових інтервалів об'єднано; застосовано словник нормалізації ASR-артефактів ($\approx$20 замін: <<сергій гс>>~$\to$~Three.js, <<афреїм>>~$\to$~A-Frame, <<Майн-Гер>>~$\to$~MindAR, <<тенсорфлоу>>~$\to$~TensorFlow.js, <<енджірок>>~$\to$~ngrok тощо); для кожного файлу обчислено частку розпізнаних україномовних сегментів як показник якості (усі файли, окрім тижня~14, мають понад 93\%).

\textbf{Метод аналізу.} Застосовано спрямований якісний контент-аналіз (\textit{directed qualitative content analysis})~\cite{Hsieh2005QCA, Mayring2014QCA}, що передбачає дедуктивне кодування за попередньо визначеними категоріями з теоретичних основ дослідження та індуктивне уточнення схеми у процесі аналізу~\cite{Kuckartz2019QTA}. Одиницею аналізу є семантичний сегмент~-- фрагмент безперервного тематичного мовлення тривалістю $\approx$10--60~с.

\textbf{Схема кодування.} Розроблено дедуктивну схему кодування за шістьма вимірами, що відповідають теоретичним конструктам розділів~1--2 дисертації (табл.~\ref{tab:coding_scheme}).

\begin{table}[!h]
\centering
\caption{Схема кодування для якісного контент-аналізу транскриптів}
\label{tab:coding_scheme}
\small
\begin{tabularx}{\linewidth}{|p{3cm}|X|p{3.9cm}|c|}
\hline
\hfil \textbf{Вимір} & \hfil \textbf{Коди} & \hfil \textbf{Теоретичне джерело} & \textbf{Підрозділ} \\
\hline
Фактори CAMIL & Присутність, Агентність, Тілесність, Когнітивне навантаження, Мотивація, Саморегуляція & Makransky \& Petersen \cite{Makransky2021CAMIL}
 & 1.1 \\
\hline
Компоненти TPACK & TK, PK, TPK, TPACK-інтегрований & Mishra \& Koehler \cite{Gurevych2015Integration} & 1.1, 2.1 \\
\hline
Педагогічні умови & Психолого-педагогічна, Методична, Технологічна, Організаційна & Авторська концепція & 2.3 \\
\hline
Етапи методики & Теоретико-аналітичний, Практико-розвивальний, Інтегративно-рефлексивний & Авторська концепція & 2.3 \\
\hline
Стратегії викладання & Моделювання, Скафолдинг, Запитування, Зворотний зв'язок, Вирішення проблем, Зв'язок з практикою & Індуктивно & --- \\
\hline
Компоненти готовності & Мотиваційно-ціннісний, Когнітивний, Діяльнісно-проєктувальний, Технологічний, Інтерактивний, Рефлексивний & Авторська модель & 2.2 \\
\hline
\end{tabularx}
\end{table}

\textbf{Процедура кодування.} Кодування здійснювалося у два етапи: (1)~автоматизоване попереднє кодування за словниками ключових слів ($\approx$150 патернів, по 8--15 на код), що забезпечило первинне маркування 73{,}3\% сегментів; (2)~інтерпретативна верифікація та уточнення автоматизованих кодів дослідником, виявлення прикладів вищого порядку інтеграції (зокрема TPK та TPACK-інтегрованого компоненту), що не піддаються ключовому пошуку, та ідентифікація індуктивних категорій. Кожний сегмент міг отримати множинні коди з різних вимірів (мультилейбельне кодування).

\textbf{Достовірність.} Забезпечення достовірності якісного аналізу~\cite{Saldana2021Coding} здійснювалося через: тривале залучення дослідника (18~тижнів систематичного спостереження за навчальним процесом як його безпосередній учасник); <<густий опис>> (\textit{thick description}) з наведенням прямих цитат із транскриптів; тріангуляцію якісних даних із кількісними результатами діагностики; прозорість кодувальної схеми (усі коди, визначення та приклади задокументовано).

\subsubsection{Зв'язок навчально-методичного забезпечення з теоретичними основами}

Формувальний етап експерименту безпосередньо реалізує структурно-функціональну модель підготовки (підрозділ~\ref{subsec:2.2}) та забезпечує виконання педагогічних умов (підрозділ~\ref{subsec:2.3}).

\textit{Реалізація блоків структурно-функціональної моделі:}
\begin{itemize}
\item \textit{мотиваційно-цільовий блок}~-- демонстрація реальних освітніх застосувань WebAR (профорієнтаційні квести \cite{Shepiliev2021Career}, віртуальна хімічна лабораторія), створення ситуацій успіху при виконанні перших завдань;
\item \textit{теоретико-методологічний блок}~-- послідовне вивчення теоретичних основ WebAR (розділи~1--3 посібника), класифікація імерсивних технологій, принципи педагогічного дизайну ІОР;
\item \textit{змістовий блок}~-- послідовне опанування технологій WebAR від основ (Three.js, WebGL) до складних застосувань (ML-інтеграція);
\item \textit{організаційно-технологічний блок}~-- систематична практична діяльність: лабораторні роботи, творчі завдання, кейс віртуальної хімічної лабораторії, підсумковий проєкт;
\item \textit{рефлексивно-оцінювальний блок}~-- захист проєктів, взаємооцінювання, есе-рефлексії, самооцінювання.
\end{itemize}

\textit{Реалізація педагогічних умов:}
\begin{itemize}
\item \textit{психолого-педагогічні}~-- підтримувальне навчальне середовище (форуми у Moodle), особистісно значущі завдання (AR-візитівка), поступове зростання складності;
\item \textit{методичні}~-- комплексний навчальний посібник (10 розділів), відеоматеріали, тести для самоперевірки, зразки виконаних завдань;
\item \textit{технологічні}~-- сучасні засоби розробки WebAR (MindAR, Three.js, TensorFlow.js, A-Frame, WebXR Device API), доступ до комп'ютерних лабораторій, ngrok для тестування, платформа Moodle;
\item \textit{організаційні}~-- чітка структура курсу (18 тижневих модулів), система оцінювання з чіткими критеріями, поєднання індивідуальної та групової роботи.
\end{itemize}

\subsubsection{Методика проведення формувального експерименту}

Методика навчання розробки ІОР у межах формувального етапу будувалася за принципом спіралеподібного нарощування компетентностей, де кожний наступний тематичний блок спирається на попередній та розширює його.

\textit{Етап~1: Фундамент (тижні~1--3).} Формування базових знань та вмінь: ознайомлення з WebAR, налаштування середовища розробки, основи Three.js (сцена, камера, рендерер, об'єкти, текстури, освітлення, анімація), знайомство з MindAR. Формування мотиваційного компоненту через демонстрацію реальних застосувань WebAR.

\textit{Етап~2: Основи відстеження (тижні~4--6).} Опанування основних типів AR-контенту: завантаження текстур та 3D-моделей (GLTF/GLB), відтворення відео та аудіо. Лабораторна робота~1 <<Візитівка у доповненій реальності>> як інтеграційне завдання.

\textit{Етап~3: Розширення можливостей (тижні~7--10).} Перехід до складнішого типу відстеження~-- Face Tracking. Опорні точки обличчя, окклюдери, маски, перемикання камер. Лабораторна робота~2 <<Примірка віртуальних аксесуарів>>.

\textit{Етап~4: Стандарт WebXR (тижні~11--14).} Найскладніший тип відстеження~-- World Tracking на основі WebXR Device API. XR-сесії, hit testing, розміщення об'єктів на реальних поверхнях, контролери.

\textit{Етап~5: Інтеграція ML (тижні~15--17).} Трикрокова методика інтеграції моделей машинного навчання у WebAR \cite{Semerikov2024WebARML_CoSinE, Semerikov2024WebAR_EduDim}:
\begin{itemize}
\item \textit{крок~1:} інтеграція стандартної моделі TensorFlow.js handpose у WebAR сцену для відстеження 21 ключової точки кисті руки та управління AR-контентом (наприклад, обчислення відстані між пальцями для керування розміром 3D-об'єкта);
\item \textit{крок~2:} навчання власної моделі класифікації зображень за допомогою Google Teachable Machine та її експорт у формат TensorFlow.js;
\item \textit{крок~3:} інтеграція власної навченої моделі у повнофункціональний WebAR додаток з відображенням результатів класифікації як AR-контенту.
\end{itemize}

\textit{Етап~6: Узагальнення та підсумковий проєкт (тиждень~18).} Знайомство з додатковими технологіями (A-Frame, model-viewer, комерційні платформи), виконання підсумкового проєкту, підсумкова атестація. Підсумковий проєкт передбачав самостійну розробку повноцінного навчального WebAR ресурсу для обраної предметної галузі з обґрунтуванням педагогічної доцільності.

Трикрокова методика є вагомим внеском дослідження, оскільки дає змогу навчити студентів використовувати ML у WebAR без необхідності глибоких знань з Data Science. Поступове ускладнення (від готових моделей через навчання власних до комплексної інтеграції) забезпечує доступність матеріалу та формує розуміння повного циклу роботи з ML у контексті імерсивних технологій.
